\documentclass[11pt, oneside]{article}
\pdfminorversion=7
\usepackage{MathematicalConceptsStyle}
\hypersetup{
    pdftitle={Mathematical Concepts for MATH261 Discrete Mathematics},
    pdfauthor={Jungwoo Kim},
    pdfsubject={Mathematical concepts for MATH261 Discrete Mathematics},
    pdfcreator={LaTeX}
}
\begin{document}
\thispagestyle{empty}
\begin{center}
    \vspace*{1.1cm}
    {\LARGE\bfseries Mathematical Concepts\par}
    \vspace{0.25cm}
    {\large for (MATH261) Discrete Mathematics\par}
    \vspace{0.75cm}
    {\normalsize Jungwoo Kim\\
    Department of Mathematics\\
    \PohangUniversityOfScienceAndTechnology\par}
\end{center}
\vspace{0.45cm}
The principal definitions and mathematical results of (MATH261) Discrete Mathematics are developed here, with proofs. Notation and mathematical writing follow \textit{(MATH000) Mathematical Language}; selected worked problems appear separately in the midterm document.

% Source: Lecture 1-2, pp. 2--4. The generalized form is named Pigeonhole Principle.
\begin{theorem}[Pigeonhole Principle]\label{thm:pigeonhole}
If $N$ objects are distributed among $k\ge1$ boxes, at least one box contains at least $\lceil N/k\rceil$ objects. In particular, distributing $k+1$ objects among $k$ boxes puts at least two objects in one box.
\end{theorem}
\begin{proof}
If every box contained fewer than $\lceil N/k\rceil$ objects, each would contain at most $\lceil N/k\rceil-1<N/k$ objects. The total would then be less than $N$, a contradiction. The particular case follows by setting $N=k+1$.
\end{proof}

% Source: Lecture 1-1, pp. 23, 28. General form used for recurrence arguments.
\begin{theorem}[Mathematical Induction]
Let $P(n)$ be a statement for each $n\in\N$.
\begin{enumerate}[label=(\roman*)]
\item If $P(1)$ holds and $P(n)$ implies $P(n+1)$ for every $n\ge1$, then $P(n)$ holds for every $n\in\N$.
\item More generally, fix $k\in\N$. If $P(1), \ldots, P(k)$ hold and $P(m-k+1), \ldots, P(m)$ together imply $P(m+1)$ for every $m\ge k$, then $P(n)$ holds for all $n\in\N$.
\end{enumerate}
\end{theorem}
\begin{proof}
In either case, suppose a counterexample exists and let $n$ be the least index for which $P(n)$ fails. In (i), the base case gives $n>1$, so $P(n-1)$ holds and the inductive step yields $P(n)$, a contradiction. In (ii), the initial cases give $n>k$, and the $k$ preceding statements are true by minimality of $n$. Applying the inductive step with $m=n-1$ again gives a contradiction.
\end{proof}

% Source: Lecture 1-2, pp. 10--21.
\begin{definition}[Binary Relation]
A \emph{binary relation} $R$ on a set $X$ is a subset of $X\times X$. We write $xRy$ to mean $(x, y)\in R$. The relation $R$ is
\begin{enumerate}[label=(\roman*)]
\item \emph{reflexive} if $xRx$ for every $x\in X$;
\item \emph{symmetric} if $xRy$ implies $yRx$ for all $x, y\in X$;
\item \emph{antisymmetric} if $xRy$ and $yRx$ together imply $x=y$;
\item \emph{transitive} if $xRy$ and $yRz$ together imply $xRz$.
\end{enumerate}
Symmetry and antisymmetry are different conditions; neither is the negation of the other.
\end{definition}

\begin{definition}[Equivalence Relation and Equivalence Class]
An \emph{equivalence relation} $\sim$ on $X$ is reflexive, symmetric, and transitive. For $x\in X$, its \emph{equivalence class} is
\[
[x]_{\sim}\coloneqq\{y\in X\mid y\sim x\}.
\]
\end{definition}

\begin{definition}[Partition]
A \emph{partition} of a set $X$ is a family $\mathcal P$ of nonempty, pairwise disjoint subsets whose union is $X$. For $X=\varnothing$, the empty family is its unique partition.
\end{definition}

\begin{theorem}[Equivalence Relations and Partitions]\label{thm:equivpartition}
The equivalence classes of an equivalence relation on $X$ form a partition of $X$. Conversely, every partition $\mathcal P$ of $X$ induces an equivalence relation given by $x\sim y$ if and only if $x$ and $y$ belong to the same block of $\mathcal P$. These constructions are inverse to each other.
\end{theorem}
\begin{proof}
Reflexivity gives $x\in[x]_{\sim}$, so the classes cover $X$ and are nonempty. If $z\in[x]_{\sim}\cap[y]_{\sim}$, then $x\sim z$ and $z\sim y$ by symmetry, hence $x\sim y$ by transitivity. Every member of $[x]_{\sim}$ is therefore equivalent to $y$, and conversely; thus $[x]_{\sim}=[y]_{\sim}$. Distinct classes are disjoint.

Conversely, every $x\in X$ lies in exactly one block of a partition. Belonging to the same block is reflexive, symmetric, and transitive. Its equivalence class at $x$ is precisely the block containing $x$, establishing both inverse assertions.
\end{proof}

\begin{definition}[Partial Order, Poset, and Total Order]
A \emph{partial order} $\preceq$ on $X$ is a reflexive, antisymmetric, and transitive relation. The pair $(X, \preceq)$ is a \emph{partially ordered set} (\emph{poset}). Two elements $x, y$ are \emph{comparable} if $x\preceq y$ or $y\preceq x$; otherwise they are \emph{incomparable}. A partial order is a \emph{total order} if every pair is comparable. We write $x\prec y$ if $x\preceq y$ and $x\ne y$.
\end{definition}

\begin{definition}[Cover Relation and Hasse Diagram]
For $x, y$ in a poset, $y$ \emph{covers} $x$, written $x\lessdot y$, if $x\prec y$ and there is no $z$ with $x\prec z\prec y$. A \emph{Hasse diagram} of a finite poset places its elements as vertices, with higher elements drawn above lower elements, and draws a segment between $x$ and $y$ precisely when $x\lessdot y$. Reflexive loops, arrowheads, and edges implied by transitivity are omitted. The order is recovered from upward paths, including paths of length zero.
\end{definition}

\begin{definition}[Extremal Elements]
Let $(X, \preceq)$ be a poset and $A\subseteq X$. An element $m\in A$ is \emph{maximal} in $A$ if no $a\in A$ satisfies $m\prec a$, and \emph{minimal} if no $a\in A$ satisfies $a\prec m$. An element $g\in A$ is \emph{greatest} in $A$ if $a\preceq g$ for every $a\in A$, and \emph{least} if $g\preceq a$ for every $a\in A$.
\end{definition}
\begin{proposition}[Greatest and Maximal Elements]
A greatest (respectively, least) element is unique, if it exists, and is maximal (respectively, minimal). A finite nonempty poset has a maximal and a minimal element, but need not have a greatest or a least element.
\end{proposition}
\begin{proof}
If $g, h$ are both greatest, then $g\preceq h$ and $h\preceq g$, so antisymmetry gives $g=h$. If $g\prec a$ for some $a$, greatestness also gives $a\preceq g$, again contradicting antisymmetry. The dual assertions follow similarly. In a finite nonempty poset, start at any element and move to a strictly greater element whenever possible. Finiteness forces termination at a maximal element. Reversing the order gives a minimal element.
\end{proof}

% Source: Lecture 1-2, pp. 22--23.
\begin{definition}[Sunflower]
A family of distinct sets $A_1, \ldots, A_r$ is a \emph{sunflower} with $r$ petals if a set $K$ exists such that $A_i\cap A_j=K$ for every $i\ne j$. The set $K$ is the \emph{core}; the sets $A_i\setminus K$ are pairwise disjoint \emph{petals}. In a $k$-uniform family every member has cardinality $k$.
\end{definition}
\Needspace{15\baselineskip}
\begin{lemma}[Sunflower Lemma]\label{lem:sunflower}
Let $k, r\in\N$ and let $\mathcal F$ be a family of distinct sets, each containing exactly $k$ elements. If
\[
|\mathcal F|>k!(r-1)^k,
\]
then $\mathcal F$ contains a sunflower with $r$ petals.
\end{lemma}
\begin{proof}
Induct on $k$. For $k=1$, more than $r-1$ distinct singleton sets have pairwise empty intersections, so any $r$ form a sunflower. The case $r=1$ is immediate for every $k$: the inequality guarantees $\mathcal F\ne\varnothing$.

Assume $k\ge2$ and $r\ge2$, and take a maximal pairwise disjoint subfamily $\{F_1, \ldots, F_t\}\subseteq\mathcal F$. If $t\ge r$, its first $r$ members form a sunflower with empty core. Otherwise $t\le r-1$. Put $U=F_1\cup\cdots\cup F_t$, so $|U|\le k(r-1)$. Maximality implies that every $F\in\mathcal F$ meets $U$. Counting pairs $(u, F)$ with $u\in U\cap F$ shows that some $u\in U$ belongs to at least
\[
\frac{|\mathcal F|}{|U|}>\frac{k!(r-1)^k}{k(r-1)}=(k-1)!(r-1)^{k-1}
\]
members. Removing $u$ from each such member produces distinct $(k-1)$-element sets. By induction these contain a sunflower with $r$ petals and core $K$; reinserting $u$ gives the required sunflower with core $K\cup\{u\}$.
\end{proof}

% Source: Lecture 2-2, pp. 2--12; MATH000 combinatorics notation.
\begin{definition}[Ordered Selections and Permutations]
For $0\le r\le n$, an \emph{ordered selection without repetition} of $r$ elements from an $n$-element set is an injective map $[r]\to[n]$, after identifying the ground set with $[n]$. The number of such selections is denoted by ${}_nP_r$. A \emph{permutation} of a set $X$ is a bijection $X\to X$; the set of permutations of $[n]$ is denoted by $S_n$.
\end{definition}
\begin{proposition}[Number of Ordered Selections]
For $0\le r\le n$,
\[
{}_nP_r=\fall nr=\frac{n!}{(n-r)!}.
\]
In particular, $|S_n|={}_nP_n=n!$, with $|S_0|=1$.
\end{proposition}
\begin{proof}
An ordered selection has $n$ choices for its first entry, $n-1$ for its second, and $n-r+1$ for its last. Multiplication yields the stated falling factorial. When $r=n$, every injection $[n]\to[n]$ is a permutation.
\end{proof}

\begin{notation}[Permutation Notation]
For $\pi\in S_n$, the symbol $\pi(i)$ denotes the image of $i$. Its \emph{two-line notation} is
\[
\pi=\begin{pmatrix}1&2&\cdots&n\\\pi(1)&\pi(2)&\cdots&\pi(n)\end{pmatrix}.
\]
The composition $\sigma\circ\pi$ sends $i$ to $\sigma(\pi(i))$, so the rightmost map acts first. The identity $\id_{[n]}$ fixes every element, and $\pi^{-1}$ denotes the inverse permutation. The cycle $(i_1\ i_2\ \cdots\ i_\ell)$ sends $i_j$ to $i_{j+1}$ for $1\le j<\ell$, sends $i_\ell$ to $i_1$, and fixes all elements not listed. A fixed point is a cycle of length one.
\end{notation}
\Needspace{13\baselineskip}
\begin{proposition}[Disjoint-Cycle Decomposition]
Every permutation of $[n]$ is a product of disjoint cycles, unique up to the order of the cycles and cyclic rotation of the entries within each cycle.
\end{proposition}
\begin{proof}
For any $i\in[n]$, the sequence $i, \pi(i), \pi^2(i), \ldots$ eventually repeats because $[n]$ is finite. Injectivity forces the first repeated element to be $i$, producing a cycle. Repeating this construction on unused elements partitions $[n]$ into disjoint cycles. The cycles are exactly the orbits of $\pi$, so their underlying sets are unique, and the entries in each cycle are determined up to cyclic rotation.
\end{proof}

\begin{definition}[Combinations]
An $r$-\emph{combination} of an $n$-element set is an $r$-element subset. Its number is denoted by $\binom nr$ (also ${}_nC_r$). We set $\binom nr=0$ when $r<0$ or $r>n$, for $n\in\Nzero$.
\end{definition}
\begin{proposition}[Number of Combinations]
For $0\le r\le n$,
\[
\binom nr=\frac{n!}{r!(n-r)!}.
\]
\end{proposition}
\begin{proof}
Each $r$-element subset has exactly $r!$ orderings. Thus the ${}_nP_r$ ordered selections are partitioned into groups of $r!$ according to their underlying subsets, and $\binom nr={}_nP_r/r!$.
\end{proof}

\begin{proposition}[Pascal's Identity]\label{prop:pascal}
For $n\ge1$ and $0\le r\le n$,
\[
\binom nr=\binom{n-1}{r-1}+\binom{n-1}{r}.
\]
\end{proposition}
\begin{proof}[Proof 1]
Fix an element $x$ of an $n$-element set. An $r$-element subset either contains $x$, leaving $r-1$ elements to be chosen from the other $n-1$, or excludes $x$, leaving $r$ elements to be chosen. These cases are disjoint and exhaustive.
\end{proof}
\begin{proof}[Proof 2]
For $1\le r\le n-1$,
\[
\begin{aligned}
\binom{n-1}{r-1}+\binom{n-1}{r}
&=\frac{(n-1)!}{(r-1)!(n-r)!}+\frac{(n-1)!}{r!(n-r-1)!}\\
&=\frac{(r+n-r)(n-1)!}{r!(n-r)!}
=\binom nr.
\end{aligned}
\]
For $r=0$ or $r=n$, the formula follows from the convention for out-of-range coefficients.
\end{proof}

\Needspace{12\baselineskip}
\begin{theorem}[Binomial Theorem]\label{thm:binomial}
For $n\in\Nzero$ and commuting variables $x, y$,
\[
(x+y)^n=\sum_{r=0}^n\binom nr x^{n-r}y^r.
\]
\end{theorem}
\begin{proof}[Proof 1]
In the product of $n$ factors $(x+y)$, obtaining $x^{n-r}y^r$ amounts to choosing the $r$ factors contributing $y$. There are $\binom nr$ such choices.
\end{proof}
\Needspace{17\baselineskip}
\begin{proof}[Proof 2]
Use induction on $n$. The identity holds at $n=0$. Assuming it at $n$ and multiplying by $x+y$,
\[
\begin{aligned}
(x+y)^{n+1}
&=\sum_{r=0}^n\binom nr x^{n+1-r}y^r
 +\sum_{r=1}^{n+1}\binom n{r-1}x^{n+1-r}y^r\\
&=x^{n+1}+\sum_{r=1}^n
\left(\binom nr+\binom n{r-1}\right)x^{n+1-r}y^r+y^{n+1}\\
&=\sum_{r=0}^{n+1}\binom{n+1}{r}x^{n+1-r}y^r,
\end{aligned}
\]
where the last step uses Pascal's identity.
\end{proof}

\begin{definition}[Multinomial Coefficient]
For $n_1, \ldots, n_k\in\Nzero$ satisfying $n_1+\cdots+n_k=n$, the \emph{multinomial coefficient} is
\[
\binom{n}{n_1, \ldots, n_k}\coloneqq\frac{n!}{n_1!\cdots n_k!}.
\]
\end{definition}
\begin{proposition}[Multinomial Counting Formula]
The number of ways to distribute $n$ distinct objects into $k$ labeled groups of prescribed sizes $n_1, \ldots, n_k$ is $\binom{n}{n_1, \ldots, n_k}$.
\end{proposition}
\begin{proof}
Order all $n$ objects, then place the first $n_1$ in group one, the next $n_2$ in group two, and so on. Each assignment occurs exactly $n_1!\cdots n_k!$ times among the $n!$ orders, corresponding to reorderings within the groups.
\end{proof}
\begin{theorem}[Multinomial Theorem]\label{thm:multinomial}
For $k\in\N$, $n\in\Nzero$, and commuting variables $x_1, \ldots, x_k$,
\[
(x_1+\cdots+x_k)^n=
\sum_{\substack{n_1, \ldots, n_k\ge0\\n_1+\cdots+n_k=n}}
\binom{n}{n_1, \ldots, n_k}\prod_{i=1}^k x_i^{n_i}.
\]
\end{theorem}
\begin{proof}[Proof 1]
Choose one term from each of the $n$ factors. The coefficient of $x_1^{n_1}\cdots x_k^{n_k}$ counts assignments of the factor positions to $k$ labeled groups of sizes $n_1, \ldots, n_k$, so it equals $\binom{n}{n_1, \ldots, n_k}$.
\end{proof}
\begin{proof}[Proof 2]
Induct on $n$. The formula is immediate for $n=0$. Assuming it for $n$, multiply by $x_1+\cdots+x_k$. For $n_1+\cdots+n_k=n+1$, the coefficient of $x_1^{n_1}\cdots x_k^{n_k}$ in the resulting product is
\[
\begin{aligned}
\sum_{\substack{1\le j\le k\\n_j\ge1}}
\binom{n}{n_1, \ldots, n_j-1, \ldots, n_k}
&=\sum_{\substack{1\le j\le k\\n_j\ge1}}
\frac{n!n_j}{n_1!\cdots n_k!}\\
&=\frac{n!}{n_1!\cdots n_k!}\sum_{j=1}^k n_j
=\binom{n+1}{n_1, \ldots, n_k}.
\end{aligned}
\]
Terms with $n_j=0$ contribute nothing, as reflected in the restricted sum. This proves the inductive step.
\end{proof}

\begin{definition}[Combinations with Repetition]
For $n\ge1$ and $r\in\Nzero$, a combination with repetition is a multiset of cardinality $r$ drawn from $n$ types. Its number is denoted by $\left(\!\binom nr\!\right)$ (also ${}_nH_r$).
\end{definition}

\begin{theorem}[Stars and Bars]\label{thm:starsbars}
For $n\in\N$ and $r\in\Nzero$, the number of solutions to
\[
x_1+\cdots+x_n=r, \qquad (x_1, \ldots, x_n)\in\Nzero^n,
\]
is $\binom{n+r-1}{n-1}=\binom{n+r-1}{r}$. It is also the number $\left(\!\binom nr\!\right)$ of $r$-combinations with repetition from $n$ types.
\end{theorem}
\begin{proof}
Represent the solution by $x_1$ stars, a bar, $x_2$ stars, and so on, ending with $x_n$ stars. This gives a string of $r$ stars and $n-1$ bars. Conversely, the numbers of stars between consecutive bars recover the $x_i$ uniquely. Choosing the positions of the $n-1$ bars among $n+r-1$ positions gives the formula. Selecting $x_i$ objects of type $i$ is the same data as the solution vector.
\end{proof}

% Source: Lecture 4-1, pp. 5--13, 14--15.
\Needspace{16\baselineskip}
\begin{theorem}[Inclusion--Exclusion Principle]\label{thm:ie}
For finite sets $A_1, \ldots, A_n$,
\[
\left|\bigcup_{i=1}^n A_i\right|=
\sum_{\varnothing\ne I\subseteq[n]}(-1)^{|I|+1}
\left|\bigcap_{i\in I}A_i\right|.
\]
Equivalently, if $A_i\subseteq U$ for a finite universe $U$, then
\[
\left|U\setminus\bigcup_{i=1}^nA_i\right|=
\sum_{I\subseteq[n]}(-1)^{|I|}\left|\bigcap_{i\in I}A_i\right|,
\qquad\bigcap_{i\in\varnothing}A_i=U.
\]
\end{theorem}
\begin{proof}[Proof 1]
Fix an element $u$, and suppose it belongs to exactly $t$ of the sets. If $t=0$, it contributes zero to both sides of the first identity. Otherwise it contributes
\[
\sum_{j=1}^t(-1)^{j+1}\binom tj=1-(1-1)^t=1
\]
to the right-hand side. Thus each element of the union is counted exactly once. Subtracting this identity from $|U|$, for the specified finite universe $U$, gives the complement formula.
\end{proof}
\begin{proof}[Proof 2]
Induct on $n$. The case $n=1$ is immediate. For $n\ge2$, the two-set identity and the inductive hypothesis give
\[
\begin{aligned}
\left|\bigcup_{i=1}^n A_i\right|
&=\left|\bigcup_{i=1}^{n-1}A_i\right|+|A_n|
-\left|\bigcup_{i=1}^{n-1}(A_i\cap A_n)\right|\\
&=\sum_{\varnothing\ne I\subseteq[n-1]}(-1)^{|I|+1}
\left|\bigcap_{i\in I}A_i\right|+|A_n|\\
&\quad-\sum_{\varnothing\ne I\subseteq[n-1]}(-1)^{|I|+1}
\left|A_n\cap\bigcap_{i\in I}A_i\right|.
\end{aligned}
\]
The first sum represents subsets of indices not containing $n$, the middle term represents $\{n\}$, and the last sum represents $I\cup\{n\}$ with the required sign $(-1)^{|I|+2}$. These cover every nonempty subset of $[n]$. The complement form follows by subtraction from $|U|$.
\end{proof}

\Needspace{12\baselineskip}
\begin{proposition}[Number of Surjections]\label{prop:surj}
Let $n, k\in\Nzero$. For $k\ge1$, the number of surjections $[n]\to[k]$ is
\[
\#\Surj([n], [k])=\sum_{j=0}^k(-1)^j\binom kj(k-j)^n.
\]
Here $0^0=1$ when an empty domain is counted. For $n=k=0$ there is one surjection $\varnothing\to\varnothing$; for $n>0$ and $k=0$ there are none.
\end{proposition}
\begin{proof}
For $k\ge1$, there are $k^n$ maps $[n]\to[k]$. Let $A_i$ be those maps whose image misses $i$. For any $J\subseteq[k]$, a map missing every element of $J$ has $(k-|J|)^n$ choices. Inclusion--exclusion counts the maps in none of the $A_i$, which are exactly the surjections.
\end{proof}

\begin{definition}[Euler's Totient Function]
For $n\in\N$, Euler's totient function is
\[
\varphi(n)\coloneqq\#\{k\in[n]\mid\gcd(k, n)=1\}.
\]
In particular, $\varphi(1)=1$.
\end{definition}
\Needspace{13\baselineskip}
\begin{theorem}[Euler's Totient Formula]\label{thm:euler}
If $n\in\N$ has distinct prime divisors $p_1, \ldots, p_t$, then
\[
\varphi(n)=n\prod_{i=1}^t\left(1-\frac1{p_i}\right)
=n\prod_{\substack{p\mid n\\p\text{ prime}}}\left(1-\frac1p\right).
\]
The empty product for $n=1$ equals $1$.
\end{theorem}
\begin{proof}
For $n\ge2$, an integer in $[n]$ fails to be coprime to $n$ precisely when it is divisible by at least one of $p_1, \ldots, p_t$. For a subset $I\subseteq[t]$, exactly $n/\displaystyle\prod_{i\in I}p_i$ integers in $[n]$ are divisible by all the primes indexed by $I$. By inclusion--exclusion, the number divisible by none is
\[
\sum_{I\subseteq[t]}(-1)^{|I|}\frac{n}{\prod_{i\in I}p_i}
=n\prod_{i=1}^t\left(1-\frac1{p_i}\right).
\]
For $n=1$, the definition and empty-product convention agree.
\end{proof}

% Source: Lecture 4-1, pp. 16--18; Lecture 4-2, pp. 20--22.
\begin{definition}[Derangements]
A \emph{derangement} of $[n]$ is a permutation $\pi\in S_n$ with no fixed point, so $\pi(i)\ne i$ for every $i\in[n]$. Denote the number of derangements by $D_n$ (also $!n$). The empty permutation is a derangement by convention, so $D_0=1$; furthermore, $D_1=0$ and $D_2=1$.
\end{definition}
\Needspace{12\baselineskip}
\begin{theorem}[Derangement Formula]\label{thm:derformula}
For every $n\in\Nzero$,
\[
D_n=n!\sum_{j=0}^n\frac{(-1)^j}{j!}.
\]
\end{theorem}
\begin{proof}
For each $i\in[n]$, let $A_i=\{\pi\in S_n\mid\pi(i)=i\}$. If $I\subseteq[n]$ has cardinality $j$, fixing all elements of $I$ leaves $(n-j)!$ permutations. Inclusion--exclusion therefore gives
\[
D_n=\sum_{j=0}^n(-1)^j\binom nj(n-j)!
=n!\sum_{j=0}^n\frac{(-1)^j}{j!}.
\]
The formula also gives $D_0=1$.
\end{proof}
\Needspace{14\baselineskip}
\begin{theorem}[Derangement Recurrences]\label{thm:derrec}
For $n\ge1$ and $n\ge2$, respectively,
\[
D_n=nD_{n-1}+(-1)^n, \qquad
D_n=(n-1)(D_{n-1}+D_{n-2}),
\]
with $D_0=1$ and $D_1=0$.
\end{theorem}
\begin{proof}[Proof 1]
By Theorem~\ref{thm:derformula},
\[
D_n-nD_{n-1}=n!\frac{(-1)^n}{n!}=(-1)^n.
\]
For $n\ge2$, using this identity at the consecutive indices $n$ and $n-1$ gives
\[
\begin{aligned}
D_n&=nD_{n-1}+(-1)^n\\
&=(n-1)D_{n-1}+\bigl((n-1)D_{n-2}+(-1)^{n-1}\bigr)+(-1)^n\\
&=(n-1)(D_{n-1}+D_{n-2}).
\end{aligned}
\]
\end{proof}
\begin{proof}[Proof 2]
We first count derangements combinatorially. Fix $j=\pi(n)\ne n$; there are $n-1$ choices. If $\pi(j)=n$, the two-cycle $(n\ j)$ can be removed, leaving a derangement of the other $n-2$ elements. Otherwise, delete $n$ from its cycle; this gives a derangement of $[n-1]$. Conversely, insert $n$ immediately before $j$ into the cycle containing $j$. This reconstructs the original permutation with $\pi(n)=j$ and $\pi(j)\ne n$. Hence, for $n\ge2$,
\[
D_n=(n-1)(D_{n-1}+D_{n-2}).
\]
For the other recurrence, rearrange this identity as
\[
D_n-nD_{n-1}=-\bigl(D_{n-1}-(n-1)D_{n-2}\bigr).
\]
Iterating the identity down to $D_1-D_0=-1$ gives
\[
D_n-nD_{n-1}=(-1)^{n-1}(D_1-D_0)=(-1)^n,
\]
as required.
\end{proof}

% Source: Lecture 4-2, pp. 2--9 (the theater problem is reserved for the midterm).
\begin{definition}[Catalan Numbers]
The \emph{Catalan numbers} are
\[
C_n\coloneqq\frac1{n+1}\binom{2n}{n}, \qquad n\in\Nzero.
\]
Thus $C_0=1$, $C_1=1$, $C_2=2$, $C_3=5$, and $C_4=14$.
\end{definition}
\begin{theorem}[Catalan Words]\label{thm:catalanwords}
For $n\in\Nzero$, the number of words $a_1\cdots a_{2n}$ with $a_i\in\{-1, 1\}$, exactly $n$ entries of each sign, and
\[
\sum_{j=1}^i a_j\ge0\qquad(1\le i\le 2n)
\]
is $C_n$.
\end{theorem}
\begin{proof}
The assertion for $n=0$ concerns the single empty word. Let $n\ge1$. There are $\binom{2n}{n}$ words with $n$ entries of each sign. Call a word \emph{bad} if some prefix sum is negative. Locate the first prefix whose sum is $-1$ and reverse the signs of all entries in that prefix. Since that prefix contains one more $-1$ than $+1$, the resulting word has $n+1$ positive entries and $n-1$ negative entries.

This map is a bijection: a word with $n+1$ positive entries first reaches prefix sum $+1$ at a unique position, and reversing the signs up to that position reconstructs a bad word. Hence the number of bad words is $\binom{2n}{n+1}$, and the desired number is
\[
\binom{2n}{n}-\binom{2n}{n+1}
=\left(1-\frac{n}{n+1}\right)\binom{2n}{n}
=C_n.
\]
\end{proof}
\begin{corollary}[Catalan Lattice Paths]\label{cor:catpaths}
The number of paths from $(0, 0)$ to $(n, n)$ consisting of $n$ right steps and $n$ upward steps that never pass above the diagonal $y=x$ is $C_n$.
\end{corollary}
\begin{proof}
Replace each right step by $+1$ and each upward step by $-1$. At every prefix, the partial sum equals the current horizontal coordinate minus the vertical coordinate. Thus remaining on or below the diagonal is equivalent to the condition in Theorem~\ref{thm:catalanwords}.
\end{proof}
\begin{remark}
The same word can be read as a balanced parenthesis string, with $+1$ an opening parenthesis and $-1$ a closing parenthesis. Nonnegative prefix sums guarantee that no closing parenthesis occurs before its matching opening parenthesis. This is a bijective interpretation, not a separate counting problem.
\end{remark}
\begin{theorem}[Catalan Recurrence]\label{thm:catalanrec}
With $C_0=1$, for every $n\ge1$,
\[
C_n=\sum_{j=0}^{n-1}C_jC_{n-1-j}.
\]
\end{theorem}
\begin{proof}
Consider a nonnegative balanced $\pm1$ word of length $2n$. Let its first return to partial sum zero occur after $2(j+1)$ entries, where $0\le j\le n-1$. Its initial block is a $+1$, followed by a nonnegative balanced word of length $2j$, followed by a $-1$. The remaining suffix is independently a nonnegative balanced word of length $2(n-1-j)$. For this fixed $j$ there are $C_jC_{n-1-j}$ choices by Theorem~\ref{thm:catalanwords}. Summing over the unique first-return index proves the recurrence.
\end{proof}

% Source: Lecture 4-2, pp. 10--14.
\Needspace{15\baselineskip}
\begin{definition}[Stirling Numbers of the First Kind]
Let $\stirone{n}{k}$ (also $c(n, k)$) be the number of permutations in $S_n$ whose disjoint-cycle decomposition has exactly $k$ cycles. This is the \emph{unsigned Stirling number of the first kind}. The \emph{signed Stirling number of the first kind} is
\[
s(n, k)\coloneqq(-1)^{n-k}\stirone{n}{k}.
\]
The initial and boundary values are $\stirone{0}{0}=s(0, 0)=1$, $\stirone{n}{0}=0$ for $n\ge1$, and $\stirone{n}{k}=0$ for $k>n$. The signed values follow from $s(n, k)=(-1)^{n-k}\stirone{n}{k}$.
\end{definition}
\begin{proposition}[First-Kind Stirling Recurrences]\label{prop:stirlingfirst}
For $n\ge1$ and $1\le k\le n$,
\[
\begin{aligned}
\stirone{n}{k}&=\stirone{n-1}{k-1}+(n-1)\stirone{n-1}{k}, \\
s(n, k)&=s(n-1, k-1)-(n-1)s(n-1, k).
\end{aligned}
\]
\end{proposition}
\begin{proof}
In a permutation of $[n]$ with $k$ cycles, either $n$ is a singleton cycle or it is not. In the first case, removing $(n)$ leaves a permutation of $[n-1]$ with $k-1$ cycles. In the second, remove $n$ from its cycle. The remaining permutation has $k$ cycles, and $n$ can be reinserted immediately after any one of the $n-1$ other elements, giving $n-1$ possibilities. This proves the first recurrence. Multiplying it by $(-1)^{n-k}$ and using $s(n, k)=(-1)^{n-k}\stirone{n}{k}$ gives the second.
\end{proof}

\begin{proposition}[First-Kind Stirling Special Values]
For $n\ge1$, $\stirone{n}{n}=1$ and $\stirone{n}{1}=(n-1)!$; for $n\ge2$, $\stirone{n}{n-1}=\binom n2$.
\end{proposition}
\begin{proof}
Only the identity permutation has $n$ cycles. An $n$-cycle can be written beginning with $1$, followed by any ordering of the remaining $n-1$ elements, yielding $(n-1)!$ permutations. To obtain $n-1$ cycles, exactly one pair of elements must form a transposition, while all other elements are fixed; there are $\binom n2$ choices.
\end{proof}

\Needspace{14\baselineskip}
\begin{proposition}[Falling-Factorial Coefficients]\label{prop:fallstir}
For $n\in\Nzero$,
\[
\fall xn=x(x-1)\cdots(x-n+1)=\sum_{k=0}^ns(n, k)x^k,
\]
with $\fall x0=1$.
\end{proposition}
\begin{proof}
The result is immediate for $n=0$. If it holds for $n-1$, then
\[
\fall xn=(x-n+1)\sum_{k=0}^{n-1}s(n-1, k)x^k.
\]
The coefficient of $x^k$ on the right is $s(n-1, k-1)-(n-1)s(n-1, k)=s(n, k)$ by Proposition~\ref{prop:stirlingfirst}.
\end{proof}

\Needspace{12\baselineskip}
\begin{definition}[Stirling Numbers of the Second Kind]
The \emph{Stirling number of the second kind} $\stirtwo{n}{k}$ (also $S(n, k)$) is the number of partitions of an $n$-element set into exactly $k$ nonempty, unlabeled blocks. We set
\[
\stirtwo{0}{0}=1, \quad \stirtwo{n}{0}=0\ (n\ge1), \quad \stirtwo{n}{k}=0\ (k>n).
\]
For $n\ge1$, $\stirtwo{n}{1}=\stirtwo{n}{n}=1$.
\end{definition}
\begin{proposition}[Second-Kind Stirling Special Value]
For $n\ge2$, $\stirtwo{n}{n-1}=\binom n2$.
\end{proposition}
\begin{proof}
A partition of $[n]$ into $n-1$ nonempty blocks consists of one two-element block and $n-2$ singleton blocks. Choosing the pair uniquely determines the partition.
\end{proof}
\begin{theorem}[Second-Kind Stirling Recurrence]\label{thm:stirlingsecond}
For $n\ge1$ and $1\le k\le n$,
\[
\stirtwo{n}{k}=\stirtwo{n-1}{k-1}+k\stirtwo{n-1}{k}.
\]
\end{theorem}
\begin{proof}
Consider the block containing $n$. If it is the singleton $\{n\}$, the other $n-1$ elements form $k-1$ blocks. Otherwise, first partition $[n-1]$ into $k$ blocks and then insert $n$ into exactly one of these $k$ blocks. The two cases are disjoint and exhaustive.
\end{proof}
\begin{theorem}[Surjections and Stirling Numbers]\label{thm:surjstirling}
For $n\in\Nzero$ and $k\in\N$,
\[
\#\Surj([n], [k])=k!\stirtwo{n}{k}, \qquad
\stirtwo{n}{k}=\frac1{k!}\sum_{j=0}^{k}(-1)^j\binom kj(k-j)^n.
\]
For $(n, k)=(0, 0)$, the corresponding cardinality identity is $1=0!\stirtwo{0}{0}$.
\end{theorem}
\begin{proof}
The nonempty fibers of a surjection $f:[n]\to[k]$ form a partition of $[n]$ into $k$ blocks. Conversely, labeling the blocks of such a partition bijectively by $[k]$ determines a surjection; there are $k!$ such labelings. Combining this counting identity with Proposition~\ref{prop:surj} gives the explicit sum.
\end{proof}
\begin{remark}
The identities in Proposition~\ref{prop:stirlingfirst} and Theorem~\ref{thm:stirlingsecond} have similar forms but count different objects: cycles of permutations versus blocks of set partitions. Their initial conditions and signs must not be interchanged.
\end{remark}

% Source: Lecture 4-2, pp. 16--19.
\begin{definition}[Bell Numbers]
The \emph{Bell number} $B_n$ counts all partitions of an $n$-element set. Hence
\[
B_n=\sum_{k=0}^n\stirtwo{n}{k}, \qquad B_0=1.
\]
The first values are $B_1=1$, $B_2=2$, $B_3=5$, $B_4=15$, and $B_5=52$.
\end{definition}
\Needspace{11\baselineskip}
\begin{theorem}[Bell Recurrence]\label{thm:bellrec}
For $n\ge1$,
\[
B_n=\sum_{k=0}^{n-1}\binom{n-1}{k}B_k.
\]
\end{theorem}
\begin{proof}
Fix one distinguished element $x$ of an $n$-element set. If exactly $n-1-k$ of the remaining elements lie in the block containing $x$, choose them in $\binom{n-1}{n-1-k}=\binom{n-1}{k}$ ways. The other $k$ elements may be partitioned in $B_k$ ways. Every partition determines a unique $k$, so summing proves the recurrence.
\end{proof}
\Needspace{12\baselineskip}
\begin{theorem}[Bell Finite-Sum Formula]\label{thm:bellsum}
For $n\in\Nzero$,
\[
B_n=\sum_{k=0}^n\frac{k^n}{k!}
\left(\sum_{j=0}^{n-k}\frac{(-1)^j}{j!}\right),
\]
where $0^0=1$ for the term occurring at $n=k=0$.
\end{theorem}
\begin{proof}
For $n\ge1$, use the explicit formula in Theorem~\ref{thm:surjstirling} and rearrange the finite sums:
\[
\begin{aligned}
B_n
&=\sum_{m=0}^n\frac1{m!}\sum_{j=0}^m(-1)^j\binom mj(m-j)^n\\
&=\sum_{m=0}^n\sum_{k=0}^m\frac{(-1)^{m-k}k^n}{(m-k)!k!}\\
&=\sum_{k=0}^n\frac{k^n}{k!}\sum_{j=0}^{n-k}\frac{(-1)^j}{j!}.
\end{aligned}
\]
For $n=0$, both sides equal $1$.
\end{proof}

% Source: Lecture 5-2 / Lecture 6-1, pp. 2--12. These files are identical.
\begin{definition}[Linear Homogeneous Recurrence Relation]
A \emph{linear homogeneous recurrence relation of degree $k$ with constant coefficients} is
\[
a_n=c_1a_{n-1}+\cdots+c_ka_{n-k}, \qquad n\ge k,
\]
where $k\ge1$, $c_1, \ldots, c_k\in\R$, and $c_k\ne0$. A \emph{solution} is a sequence satisfying the relation at every indicated index. We may also consider complex-valued solutions, with the same coefficients.
\end{definition}
\begin{proposition}[Linearity of Homogeneous Solutions]\label{prop:linrec}
If $(a_n)$ and $(b_n)$ solve the same linear homogeneous recurrence, then $(\alpha a_n+\beta b_n)$ also solves it for any scalars $\alpha, \beta$ in the underlying field.
\end{proposition}
\begin{proof}
For each $n\ge k$, multiply the recurrence for $a_n$ by $\alpha$, the recurrence for $b_n$ by $\beta$, and add. The result is
\[
\alpha a_n+\beta b_n=\sum_{j=1}^kc_j(\alpha a_{n-j}+\beta b_{n-j}).
\]
\end{proof}
\begin{proposition}[Existence and Uniqueness from Initial Values]\label{prop:unique}
For any prescribed $a_0, \ldots, a_{k-1}$, a recurrence of degree $k$ as above determines exactly one sequence $(a_n)_{n\ge0}$.
\end{proposition}
\begin{proof}
For existence, compute $a_k$ from $a_0, \ldots, a_{k-1}$, and then calculate each subsequent term from the preceding $k$ values. For uniqueness, two solutions with the same initial values agree at indices $0, \ldots, k-1$; if their preceding $k$ terms agree, the recurrence forces their next terms to agree. Induction completes the argument.
\end{proof}

\Needspace{11\baselineskip}
\begin{definition}[Characteristic Equation]
The \emph{characteristic polynomial} of the recurrence
$a_n=\displaystyle\sum_{j=1}^kc_ja_{n-j}$ is
\[
p(r)\coloneqq r^k-c_1r^{k-1}-\cdots-c_{k-1}r-c_k.
\]
The equation $p(r)=0$ is its \emph{characteristic equation}.
\end{definition}
\begin{proposition}[Characteristic-Root Solutions]
For $r\in\C\setminus\{0\}$, the sequence $(r^n)_{n\ge0}$ solves the homogeneous recurrence if and only if $p(r)=0$.
\end{proposition}
\begin{proof}
Substitution gives
\[
r^n-\sum_{j=1}^kc_jr^{n-j}=r^{n-k}p(r).
\]
Because $r\ne0$, the recurrence holds for all $n\ge k$ exactly when $p(r)=0$.
\end{proof}

\begin{theorem}[Second-Order Recurrences with Distinct Roots]\label{thm:distinct2}
Suppose $c_2\ne0$ and $r^2-c_1r-c_2$ has two distinct roots $r_1, r_2\in\C$. Then the solutions of
\[
a_n=c_1a_{n-1}+c_2a_{n-2}\qquad(n\ge2)
\]
are precisely
\[
a_n=\alpha r_1^n+\beta r_2^n, \qquad\alpha, \beta\in\C.
\]
For real initial values, the resulting sequence is real; nonreal roots occur in conjugate pairs.
\end{theorem}
\begin{proof}
Each $(r_i^n)$ is a solution by the characteristic equation, and linearity gives every displayed combination. Conversely, given $a_0, a_1$, the equations
\[
\alpha+\beta=a_0, \qquad\alpha r_1+\beta r_2=a_1
\]
have the unique solution
\[
\alpha=\frac{a_1-a_0r_2}{r_1-r_2}, \qquad
\beta=\frac{a_0r_1-a_1}{r_1-r_2}.
\]
The resulting combination agrees with the prescribed initial values, hence with the original solution by Proposition~\ref{prop:unique}.
\end{proof}

\begin{theorem}[Second-Order Recurrences with a Repeated Root]\label{thm:repeated2}
Suppose $c_2\ne0$ and $r^2-c_1r-c_2=(r-r_0)^2$. Then all solutions of the associated second-order recurrence are
\[
a_n=(\alpha+\beta n)r_0^n, \qquad\alpha, \beta\in\C.
\]
\end{theorem}
\begin{proof}
Since $r_0\ne0$, coefficient comparison gives $c_1=2r_0$ and $c_2=-r_0^2$. Besides $r_0^n$, the sequence $nr_0^n$ is a solution, because
\[
2r_0(n-1)r_0^{n-1}-r_0^2(n-2)r_0^{n-2}
=nr_0^n.
\]
By linearity, every displayed combination is a solution. To match arbitrary initial values $a_0, a_1$, choose $\alpha=a_0$ and $\beta=a_1/r_0-a_0$. Uniqueness of the recurrence then establishes that every solution has this form.
\end{proof}

\Needspace{20\baselineskip}
\begin{theorem}[General Solution of a Constant-Coefficient Homogeneous Recurrence]\label{thm:generalrec}
Let $p(r)=r^k-c_1r^{k-1}-\cdots-c_k$, where $c_k\ne0$. Suppose its distinct complex roots are $r_1, \ldots, r_t$ with multiplicities $m_1, \ldots, m_t$, so $m_1+\cdots+m_t=k$. The solutions of the recurrence are exactly
\[
a_n=\sum_{i=1}^t P_i(n)r_i^n,
\qquad\deg P_i<m_i,
\]
where each $P_i$ is a polynomial over $\C$. The $k$ polynomial coefficients are uniquely determined by $k$ initial values. For real recurrence coefficients and real initial values, the expression can be chosen in conjugate pairs to yield a real-valued sequence.
\end{theorem}
\begin{proof}
Let $E$ be the forward shift operator on sequences, $(Ea)_n=a_{n+1}$, and let $\Delta P(n)=P(n+1)-P(n)$. The recurrence is equivalent to $p(E)a=0$ (with the equation indexed starting at $n=0$). For $r\ne0$,
\[
(E-r)\bigl(P(n)r^n\bigr)=r^{n+1}\Delta P(n).
\]
Consequently, $(E-r)^m(P(n)r^n)=r^{n+m}\Delta^mP(n)=0$ whenever $\deg P<m$. Since the factors of $p(E)$ commute, every sequence in the displayed form is a solution.

These $k$ basic solutions are linearly independent. Indeed, suppose
\[
\sum_{i=1}^tP_i(n)r_i^n=0\qquad(n\ge0),
\]
where $\deg P_i<m_i$. Fix $i$ and apply
\[
\prod_{j\ne i}(E-r_j)^{m_j}.
\]
All terms except the $i$th disappear. On sequences $P(n)r_i^n$, each factor $E-r_j$, for $j\ne i$, transforms $P$ into
\[
(r_i-r_j)P+r_i\Delta P.
\]
On the vector space of polynomials of degree less than $m_i$, this map is triangular in the monomial basis with nonzero diagonal $r_i-r_j$, hence invertible. Therefore the surviving term can vanish only if $P_i=0$. As $i$ was arbitrary, all $P_i$ vanish.

Thus the $k$ basic solutions have independent initial-value vectors $(a_0, \ldots, a_{k-1})\in\C^k$: otherwise their linear combination would have $k$ zero initial terms and hence be identically zero by Proposition~\ref{prop:unique}. These $k$ vectors form a basis of $\C^k$, so every choice of initial values determines exactly one linear combination of the proposed form.
\end{proof}

% Source: Lecture 5-2 / Lecture 6-1, pp. 14--20.
\begin{definition}[Linear Nonhomogeneous Recurrence Relation]
A \emph{linear nonhomogeneous recurrence relation of degree $k$ with constant coefficients} is
\[
a_n=c_1a_{n-1}+\cdots+c_ka_{n-k}+F(n), \qquad n\ge k,
\]
where $c_1, \ldots, c_k$ are constants, $c_k\ne0$, and $F$ is a prescribed function. The recurrence obtained by setting $F=0$ is the \emph{associated homogeneous recurrence}. A \emph{particular solution} is any one sequence satisfying the nonhomogeneous recurrence.
\end{definition}
\Needspace{12\baselineskip}
\begin{theorem}[Homogeneous Plus Particular Solution]\label{thm:particular}
If $a_n^{(p)}$ is a particular solution to a linear nonhomogeneous recurrence, then every solution is of the form
\[
a_n=a_n^{(p)}+a_n^{(h)},
\]
where $a_n^{(h)}$ is an arbitrary solution of the associated homogeneous recurrence.
\end{theorem}
\begin{proof}
Subtract the recurrence satisfied by $a_n^{(p)}$ from the recurrence satisfied by any other solution $a_n$. Their difference $b_n=a_n-a_n^{(p)}$ satisfies $b_n=\displaystyle\sum_{j=1}^kc_jb_{n-j}$ because the forcing terms cancel. Conversely, adding a homogeneous solution to the particular solution preserves the forcing term $F(n)$, yielding another solution.
\end{proof}

\begin{proposition}[Method of Undetermined Coefficients]\label{prop:undetermined}
Let $p(r)$ be the characteristic polynomial of a linear recurrence with constant coefficients, and suppose its forcing term is $F(n)=P_d(n)\lambda^n$, where $\lambda\in\C\setminus\{0\}$ and $P_d$ is a polynomial of degree at most $d$. If $\lambda$ is a root of $p$ of multiplicity $m$ (and $m=0$ if it is not a root), a particular solution may be sought in the form
\[
a_n^{(p)}=n^mQ_d(n)\lambda^n,
\]
where $Q_d$ is a polynomial of degree at most $d$ with undetermined coefficients. For a sum of such forcing terms, particular solutions may be added.
\end{proposition}
\begin{proof}
Let $E$ denote the forward shift, $(Ea)_n=a_{n+1}$, and let $\Delta Q(n)=Q(n+1)-Q(n)$. With the recurrence indexed from $n=k$, its operator form is
\[
(p(E)a)_n=F(n+k)\qquad(n\ge0).
\]
For $a_n=\lambda^nQ(n)$, direct calculation gives
\[
(p(E)a)_n=\lambda^np\bigl(\lambda(1+\Delta)\bigr)Q(n).
\]
If $\lambda$ has multiplicity $m$, write $p(\lambda(1+z))=z^mh(z)$ with $h(0)\ne0$. The operator $h(\Delta)$ is invertible on polynomials of bounded degree: its matrix in the monomial basis is triangular with nonzero diagonal $h(0)$. Furthermore, $\Delta^m$ maps polynomials $n^mQ_d(n)$ with $\deg Q_d\le d$ bijectively onto polynomials of degree at most $d$, since the image of $n^{m+j}$ has leading term $(m+j)!n^j/j!$. Consequently the coefficients of $Q_d$ can be uniquely chosen to match $\lambda^kP_d(n+k)$ on the right. Superposition follows from linearity.
\end{proof}

\begin{center}
\Needspace{12\baselineskip}
\small
\renewcommand{\arraystretch}{1.22}
\begin{tabular}{@{}p{0.35\linewidth}p{0.53\linewidth}@{}}
\toprule
Forcing term $F(n)$ & Particular-solution trial \\
\midrule
$P_d(n)$ & $n^mQ_d(n)$, where $m$ is the multiplicity of $1$ as a characteristic root \\
$\lambda^n$ ($\lambda\ne0$) & $n^mC\lambda^n$, where $m$ is the multiplicity of $\lambda$ \\
$P_d(n)\lambda^n$ ($\lambda\ne0$) & $n^mQ_d(n)\lambda^n$ \\
Sum of terms of these forms & Add the corresponding particular-solution trials \\
\bottomrule
\end{tabular}
\end{center}
\begin{remark}
The method specifies a finite-dimensional trial family, not a guess without justification. After finding its coefficients, add the general homogeneous solution from Theorem~\ref{thm:generalrec} and use the initial conditions. Here $P_d$ is prescribed, whereas $Q_d$ has coefficients to be determined. A constant or polynomial forcing term corresponds to $\lambda=1$.
\end{remark}

% Source: Lecture 5-2 / Lecture 6-1, pp. 21--26.
\begin{definition}[Linear Recurrences with Variable Coefficients]
A \emph{linear recurrence with variable coefficients} has the form
\[
a_n=\sum_{j=1}^kc_j(n)a_{n-j}+F(n),
\]
where $n\ge k$, $k\ge1$, and the coefficients $c_j(n)$ may depend on $n$. Unlike the constant-coefficient case, its solutions are not generally determined by a single characteristic polynomial.
\end{definition}
\Needspace{13\baselineskip}
\begin{proposition}[Reduction by Rescaling]
Suppose a recurrence can be written as
\[
f(n)a_n=\sum_{j=1}^kc_jf(n-j)a_{n-j}+F(n),
\]
where $c_1, \ldots, c_k$ are constants. Setting $b_n=f(n)a_n$ gives
\[
b_n=\sum_{j=1}^kc_jb_{n-j}+F(n).
\]
If $f(n)\ne0$, one recovers $a_n=b_n/f(n)$; when $f$ vanishes at an index, that index must be handled separately.
\end{proposition}
\begin{proof}
Substitute $b_{n-j}=f(n-j)a_{n-j}$ for $0\le j\le k$. The resulting relation has constant coefficients.
\end{proof}
\Needspace{22\baselineskip}
\begin{proposition}[First-Order Variable-Coefficient Recurrences]\label{prop:varfirst}
Let
\[
a_{n+1}=f(n)a_n+g(n), \qquad n\in\Nzero,
\]
where $a_0$ is given. With the convention that empty products equal $1$, the unique solution is
\[
\boxed{\displaystyle
a_n=a_0\prod_{j=0}^{n-1}f(j)
+\sum_{m=0}^{n-1}g(m)\prod_{j=m+1}^{n-1}f(j).}
\]
This formula remains valid even if some $f(j)=0$. If $f(j)\ne0$ for $0\le j<n$, it can also be written as
\[
a_n=\left(\prod_{j=0}^{n-1}f(j)\right)
\left(a_0+\sum_{m=0}^{n-1}\frac{g(m)}{\prod_{j=0}^{m}f(j)}\right).
\]
\end{proposition}
\begin{proof}
For $n=0$, the first formula gives $a_0$. Suppose it holds at $n$. Multiplying by $f(n)$ and adding $g(n)$ yields
\[
\begin{aligned}
a_{n+1}
&=a_0\prod_{j=0}^{n}f(j)
+\sum_{m=0}^{n-1}g(m)\prod_{j=m+1}^{n}f(j)+g(n)\\
&=a_0\prod_{j=0}^{n}f(j)
+\sum_{m=0}^{n}g(m)\prod_{j=m+1}^{n}f(j),
\end{aligned}
\]
which is the formula at $n+1$. Induction proves the first expression, and forward recursion gives uniqueness. When all factors are nonzero, factoring out the cumulative product gives the second expression. Equivalently, one may divide the recurrence by that product and telescope.
\end{proof}
\begin{remark}
The product--sum formula is preferable to a quotient formula when some coefficient vanishes. In particular, a recurrence with $f(0)=0$ must not be solved by dividing by $f(0)$; one may use the first formula or begin the quotient argument after the exceptional index.
\end{remark}

% Source: Lecture 7-1, pp. 2--3.
\begin{definition}[Ordinary Generating Function]
For a sequence $(a_n)_{n\ge0}$ over $\R$ or $\C$, its \emph{ordinary generating function} (OGF) is the formal power series
\[
A(x)\coloneqq\sum_{n=0}^{\infty}a_nx^n.
\]
The notation $[x^n]A(x)$ denotes the coefficient $a_n$. Unless convergence is explicitly asserted, an identity of generating functions is an identity of formal power series, and does not require assigning a numerical value to $x$.
\end{definition}
\begin{proposition}[Ordinary Generating Function Operations]\label{prop:ogfops}
Let $A(x)=\displaystyle\sum_{n\ge0}a_nx^n$ and $B(x)=\displaystyle\sum_{n\ge0}b_nx^n$. Then
\[
\begin{aligned}
A'(x)&=\sum_{n\ge0}(n+1)a_{n+1}x^n, \\
A(x)+B(x)&=\sum_{n\ge0}(a_n+b_n)x^n, \\
A(x)B(x)&=\sum_{n\ge0}\left(\sum_{j=0}^na_jb_{n-j}\right)x^n.
\end{aligned}
\]
In particular, $[x^n](AB)=\displaystyle\sum_{j=0}^na_jb_{n-j}$ is the \emph{Cauchy product} formula.
\end{proposition}
\begin{proof}
Formal differentiation gives $\displaystyle\sum_{n\ge1}na_nx^{n-1}$, and reindexing yields the first equality. Addition is coefficientwise. In the product, a monomial $x^n$ arises precisely by multiplying $a_jx^j$ by $b_{n-j}x^{n-j}$ for $j=0, \ldots, n$, proving the third equality. Each coefficient of a product involves only finitely many summands, so no convergence assumption is needed.
\end{proof}

% Source: Lecture 7-1, pp. 5--7.
\begin{definition}[Generalized Binomial Coefficient]
For $\alpha\in\R$ and $n\in\Nzero$, define
\[
\binom\alpha0\coloneqq1, \qquad
\binom\alpha n\coloneqq\frac{\alpha(\alpha-1)\cdots(\alpha-n+1)}{n!}
=\frac{\fall\alpha n}{n!}\quad(n\ge1).
\]
When $\alpha$ is a nonnegative integer, this agrees with the usual binomial coefficient and vanishes for $n>\alpha$.
\end{definition}
\Needspace{14\baselineskip}
\begin{theorem}[Generalized Binomial Theorem]\label{thm:genbinomial}
For every $\alpha\in\R$,
\[
(1+x)^\alpha=\sum_{n=0}^{\infty}\binom\alpha n x^n
\qquad(|x|<1).
\]
The same coefficients define the corresponding formal binomial expansion. For nonnegative integer $\alpha$, the series terminates and is the ordinary Binomial Theorem.
\end{theorem}
\begin{proof}[Proof 1]
Set $F(x)=\displaystyle\sum_{n\ge0}\binom\alpha n x^n$. Its coefficients satisfy
\[
(n+1)\binom\alpha{n+1}=(\alpha-n)\binom\alpha n.
\]
The ratio test gives absolute convergence for $|x|<1$ unless the sum terminates. Termwise differentiation and comparison of coefficients then yield
\[
(1+x)F'(x)=\alpha F(x), \qquad F(0)=1.
\]
It follows that $\frac{d}{dx}((1+x)^{-\alpha}F(x))=0$ for $|x|<1$, so $F(x)=(1+x)^\alpha$. The coefficient recurrence also determines the formal expansion without requiring convergence.
\end{proof}
\begin{proof}[Proof 2]
The Maclaurin coefficients of $f(x)=(1+x)^\alpha$ satisfy
\[
f^{(n)}(x)=\fall\alpha n (1+x)^{\alpha-n},
\qquad\frac{f^{(n)}(0)}{n!}=\binom\alpha n.
\]
The function $(1+z)^\alpha=\exp(\alpha\log(1+z))$ is analytic on the complex disk $|z|<1$, with the logarithm branch satisfying $\log 1=0$. Hence its Maclaurin series represents the function throughout that disk. Restricting to real $x$ gives the claimed expansion.
\end{proof}
\begin{remark}
For $m\in\N$, negative exponents give
$\binom{-m}{n}=(-1)^n\binom{m+n-1}{n}$.
The corresponding explicit expansions of $(1+x)^{-m}$ and $(1-x)^{-m}$ may be obtained from Theorem~\ref{thm:genbinomial}; their derivation as worked examples belongs in the midterm document.
\end{remark}

% Source: Lecture 7-1, pp. 18--19.
\Needspace{25\baselineskip}
\begin{definition}[Exponential Generating Function]
For a sequence $(a_n)_{n\ge0}$ over $\R$ or $\C$, its \emph{exponential generating function} (EGF) is
\[
\widehat A(x)\coloneqq\sum_{n=0}^{\infty}a_n\frac{x^n}{n!}.
\]
Its coefficients recover the original sequence by $a_n=n![x^n]\widehat A(x)$. As for OGFs, formal power-series manipulations do not require convergence.
\end{definition}
\Needspace{13\baselineskip}
\begin{proposition}[Exponential Generating Function Operations]\label{prop:egfops}
Let $\widehat A(x)=\displaystyle\sum_{n\ge0}a_nx^n/n!$ and $\widehat B(x)=\displaystyle\sum_{n\ge0}b_nx^n/n!$. Then
\[
\begin{aligned}
\widehat A'(x)&=\sum_{n\ge0}a_{n+1}\frac{x^n}{n!}, \\
\widehat A(x)+\widehat B(x)&=\sum_{n\ge0}(a_n+b_n)\frac{x^n}{n!}, \\
\widehat A(x)\widehat B(x)&=\sum_{n\ge0}
\left(\sum_{j=0}^n\binom nj a_jb_{n-j}\right)\frac{x^n}{n!}.
\end{aligned}
\]
\end{proposition}
\begin{proof}
Differentiation of $a_nx^n/n!$ gives $a_nx^{n-1}/(n-1)!$ for $n\ge1$, yielding the first identity after reindexing. The second follows coefficientwise. For the third, the coefficient of $x^n$ in the product is
\[
\sum_{j=0}^n\frac{a_j}{j!}\frac{b_{n-j}}{(n-j)!}
=\frac1{n!}\sum_{j=0}^n\binom nj a_jb_{n-j},
\]
which gives the asserted EGF coefficient formula.
\end{proof}

\end{document}
